\documentclass[10pt,a4paper]{amsart}
\usepackage[margin=19mm]{geometry}
\usepackage{amsmath,amssymb,mathtools}
\usepackage[scr=rsfs,cal=euler]{mathalfa}
\usepackage{xfrac}
\usepackage[unicode,hidelinks,bookmarks=false]{hyperref}
\newcommand{\ie}{i.e.\ }
\newcommand{\cf}{cf.\ }
\newcommand{\resp}{respectively\ }
\newcommand{\Subsection}{\subsection}
\providecommand{\nobreakdash}{\nobreak\hbox{-}}
\setlength{\emergencystretch}{2em}
\multlinegap0pt
\usepackage{bm}
\usepackage{bbm}
\usepackage{dsfont}
\usepackage{xspace}
\usepackage{cancel}
\usepackage{mathtools}
\usepackage{amsthm}
\makeatletter
\@ifundefined{thm}{\newtheorem{thm}{Thm}}{}
\@ifundefined{lem}{\newtheorem{lem}[thm]{Lemma}}{}
\makeatother
\newcommand{\bb}[1]{{\mathbb{#1}}} 
\title[Existence and uniqueness of Remotely Almost Periodic solutio]{Existence and uniqueness of Remotely Almost Periodic solutions of differential equations and applications}
\author{Diego Jaure, Christopher Maulen}
\date{Source: arXiv:2509.11069v1 (CC BY 4.0)}
\begin{document}
\maketitle

\noindent\emph{Source and licence.} Existence and uniqueness of Remotely Almost Periodic solutions of differential equations and applications. Version arXiv:2509.11069v1 (CC BY 4.0), distributed under the Creative Commons Attribution 4.0 International licence, as recorded in the arXiv licence metadata for this exact version and shown on the abstract page https://arxiv.org/abs/2509.11069v1. Full rights notice: licenses/arxiv-2509.11069-CC-BY-4.0.txt.\par\smallskip

\noindent\textit{Editorial scope.} Counted unit: the Lem whose source label is \texttt{lem\_transformacion} in the article (source0.tex lines 760--782, source label \texttt{lem\_transformacion}), delivered together with the article's own complete original proof of it (source0.tex lines 784--937, one \texttt{proof} environment of 154 lines). No mathematics is written, repaired or re-derived: statement and proof are the article's own text line for line. Required context retained uncounted: averaging\_inicial (lines 639--641); promedio (lines 652--654); lema\_tecnico\_desigualdad (lines 684--704). Every definition, display and label of those lines is retained unchanged. Omitted from this package and not redistributed: 1--638, 642--651, 655--683, 705--759, 783--783, 938--1425. Nothing inside a retained excerpt is omitted or altered: every retained body is delivered exactly as the article's lines are written, so no omission receipt of any kind accompanies this package. Citations: the retained excerpts carry no citation command at all, so no bibliography entry of the article is transcribed and no reference is redistributed. Reference adaptation: none was needed and none was made; every reference inside the delivered unit resolves inside it, so no unresolved reference or citation is produced. Printed slips kept literal: no printed word, symbol, hypothesis or number of the counted statement or of its proof was altered. Layout and class: the article's own class is \texttt{amsart} with packages including fontenc, inputenc, geometry, verbatim, amsthm, amssymb, setspace, enumerate, fancyhdr, enumitem, hyperref, babel; the wrapper is the lane's pinned \texttt{amsart} editorial wrapper at 10pt on A4, which restates the theorem-like environments the retained text needs and adds the article's own macro definitions that the retained text needs (\texttt{\textbackslash bb}), with extra wrapper packages (bm, bbm, dsfont, xspace, cancel, mathtools). The delivered numbering is the unit's own. Assets: no retained span contains a figure environment, a graphics-inclusion command, a PDF-page inclusion, a table or a TikZ picture; no image file is copied into this package, the package declares no asset dependency and no image file is redistributed with it. Licence: version arXiv:2509.11069v1 of this article is distributed under the Creative Commons Attribution 4.0 International licence, as recorded in the arXiv licence metadata for this exact version and shown on the abstract page https://arxiv.org/abs/2509.11069v1. A full rights notice is delivered at licenses/arxiv-2509.11069-CC-BY-4.0.txt. Characters: every retained excerpt is pure ASCII, so the delivered engine, XeLaTeX with its default Latin Modern text font, reports no missing character.
\begin{eqnarray}  
\frac{dx}{dt}=\nu f_{\nu}(t,x),\label{averaging_inicial}  
\end{eqnarray}  

\begin{eqnarray}
f_0(x)=\lim_{T\rightarrow \infty} \frac{1}{2T}\int_{-T}^{T} f_0(t,x)dt, x\in B_{r}(0). \label{promedio}
\end{eqnarray}

\begin{lem}\label{lema_tecnico_desigualdad}
Let $f\in C(\bb{R}\times\Omega,\bb{R})$, we define $F:\bb{R}\times\Omega\times (0,\infty)\rightarrow \bb{C}$ by
\begin{eqnarray}
F(t,x,\nu)=\int_{-\infty}^{t} e^{-\nu(t-s)}f(s,x)ds.\label{F_exp_nu}
\end{eqnarray}
Moreover let
\begin{eqnarray}
h(u,x)&=&\sup_{u\in\bb{R}}\left\vert \frac{1}{u}\int_{0}^{u}f(s-t,x)dt \right\vert \label{h_ux}\\
\xi(x,r)&=&r^{2}\int_{0}^{\infty} h(u,x) ue^{-ru}du.\label{xi_xr}
\end{eqnarray}
Then we have
\begin{enumerate}
\item $\vert F(t,x,\nu)\vert\leq \nu^{-1}\xi(x,\nu)$;

\item  $\vert \frac{\partial F}{\partial t}(t,x,\nu)-f(t,x)\vert\leq \xi(x,\nu)$

\item If $f$ is ergodic, then $F$ is also ergodic, with $M(F)(x)=\nu^{-1}M(f)(x)$, for fixed $\nu>0$.

\item If $f\in RAP(\bb{R}\times\Omega,\bb{R})$ then $F\in RAP(\bb{R}\times\Omega,\bb{R})$ for $\nu>0.$
\end{enumerate}
\end{lem}

\begin{lem}\label{lem_transformacion}
If $f_\nu$ satisfies the conditions mentioned in the first paragraph of this section.  
Let $f_0$ be as in \eqref{promedio}. Then, for all $r<r_0$, there exists $\nu_0>0$ and a continuous function $U$ on $\mathbb{R}\times B_{r}(0)\times (0,\infty)$ such that
\begin{enumerate}
\item For each $\nu\in(0,\infty)$, $U\in RAP(\mathbb{R}\times B_{r}(0),\mathbb{R}^n)$, and it is ergodic, that is, its average exists.

\item $\frac{\partial U}{\partial t}$ is continuous on $\mathbb{R}\times \mathbb{R}^n\times (0,\infty)$, and derivatives of arbitrary order with respect to $x\in\mathbb{R}^n$ are continuous for each $\nu\in (0,\infty)$.  
$\frac{\partial U}{\partial t}$ and these derivatives are remotely almost-periodic and ergodic.

\item Let $G(t,x,\nu)=\frac{\partial U}{\partial t}(t,x,\nu)-f_{\nu}(t,x)+f_{0}(x)$.  
Then all the functions $\nu U,\ \nu \frac{\partial U}{\partial x},\ G$, and $\frac{\partial G}{\partial x}$ tend to zero as $\nu\rightarrow 0$, uniformly on $\mathbb{R}\times B_{r}(0)$.

\item The change of variable  
\begin{equation}
x=y+\nu U(t,y,\nu) \quad \text{with } (t,y,\nu)\in\mathbb{R}\times B_{r}(0)\times [0,\nu_0] \label{cambio_variable}
\end{equation}  
is invertible and transforms \eqref{averaging_inicial} into  
\begin{eqnarray}
\frac{\partial y}{\partial t}=\nu f_{0}(y)+\nu F_{\nu}(t,y)\label{reducida}
\end{eqnarray}  
where $F_{\nu}$ satisfies the same properties as $f_{\nu}$ on $\mathbb{R}\times B_{r}(0)\times[0,\nu_{0}]$, and additionally, $F_{0}(t,y)=0$ for all $(t,y)\in\mathbb{R}\times B_{r}(0)$.
\end{enumerate}
\end{lem}

\begin{proof}
Let $H:\bb{R}\times B_{r}(0)\rightarrow \bb{C}^n$ by
$$H(t,x)=f_0 (t,x)-f_0 (x)$$
then $H\in RAP(\bb{R}\times B_{r}(0),\bb{R}^n)$. As $f$ is ergodic, it follows from \eqref{promedio} that
$$\lim_{T\rightarrow \infty}\frac{1}{2T}\int_{-T}^{T} H(t+s,x)dt =0$$
uniformly with respect to $x\in B_{r}(0)$ and $s\in \bb{R}$. Therefore $H$ is ergodic and $\mathcal{M}(H_x)=0$. 
Let us define the functions $h:\bb{R}\times B_{r}(0)\rightarrow \mathbb{R}$ and $\xi:\bb{R}\times B_{r}(0)\rightarrow \mathbb{R}$ given by
\begin{eqnarray*}
h(t,x)&:=&\sup_{s\in\bb{R}}\left\vert \frac{1}{t}\int_{0}^{t}H(s-u,x)du\right\vert \,\, x\in B_{r}(0)\\
\xi(t,x)&:=&t^2\int_{0}^{\infty} e^{-tu}u h(u,x)du
\end{eqnarray*}
for every $(t,x)\in\bb{R}\times B_{r}(0)$. It is straightforward that $\xi(x,t)\rightarrow 0$ when $t\rightarrow 0$ uniformly with respect to $x\in B_{r}(0)$. Thus, we define $\xi(x,0):=0$ for every $x\in B_{r}(0)$.

Consider $\overline{H}:\bb{R}\times B_{r}(0) \times (0,\infty)\rightarrow\bb{C}^n$ the bounded solution of the differential equation
\begin{eqnarray*}
\frac{\partial}{\partial t} \overline{H}(t,x,\nu)-H(t,x)=-\nu \overline{H}(t,x,\nu),
\end{eqnarray*}
given by
\begin{eqnarray*}
\overline{H}(t,x,\nu)=\int_{-\infty}^{t}e^{-\nu (t-s)}H(s,x)ds.
\end{eqnarray*}
By Lemma \ref{lema_tecnico_desigualdad}, we have that
\begin{eqnarray}
\vert \overline{H}(t,x,\nu) \vert\leq \nu^{-1}\xi(\nu,x),\ t\in\bb{R}. \label{desigualdad_overlineH}
\end{eqnarray}
Moreover, $\overline{H} \in RAP(\mathbb{R}\times B_{r}(0),\mathbb{R}^n)$ and it is ergodic with $\mathcal{M}(H_x)=0$ for fixed $\nu\in(0,\infty)$.  
It is easy to see that $\frac{\partial \overline{H}}{\partial t} \in RAP(\mathbb{R}\times B_{r}(0),\mathbb{R}^n)$ and is ergodic for $\nu\in (0,\infty)$.  
By \eqref{desigualdad_overlineH}
\begin{eqnarray*}
\left\vert \frac{\partial}{\partial t} \overline{H}(t,x,\nu)-H(t,x) \right\vert \leq \xi(\nu,x), \ t\in \bb{R}.
\end{eqnarray*}
For fixed $a>0$ and some integer $q\geq 1$, we define $\Delta_a$ in $\mathbb{C}^n$ such that
\begin{eqnarray*}
\Delta_{a}(x)=\begin{cases}
d_{a}(1-a^{-2}\left|x\right|^{2})^{2q} & \mbox{si }\left|x\right|\leq a,\\
0 & \mbox{si }\left|x\right|>a,
\end{cases}
\end{eqnarray*}
where the constant $d_a$ is determined by
\begin{eqnarray*}
\int_{B_{a}} \Delta_a (x)dx=1.
\end{eqnarray*}
Let us define now the function $U:\bb{R}\times\bb{C}^n\times(0,\infty)\rightarrow\bb{C}^n$ by
\begin{eqnarray*}
U(t,x,\nu)=\int_{B_{a}} \Delta_a(x-y)\overline{H}(t,y,\nu)dy.
\end{eqnarray*}
which is a continuous functions, also %theo 1.7
$U \in RAP(\mathbb{R}\times B_{a})$ for $\nu\in(0,\infty)$, and $U$ is ergodic since $\overline{H}$ is. Therefore, condition (1) is satisfied.

To prove (2). The function $\Delta_a(x - y)$ has continuous partial derivatives of order greater than $2q - 1$ with respect to $x$, which are bounded, in norm, by a function $A(a)$ (the integration area), where $A$ is continuous on $(0,\infty)$.  
From \eqref{desigualdad_overlineH}, it follows that the function $U$ has partial derivatives with respect to $x$ of order greater than $2q - 1$, which are bounded by $A(a)\xi(y,\nu)\nu^{-1}$, for $y\in B_r(0)$.  
Since $q$ is an arbitrary integer, the number of derivatives with respect to $x$ can be as large as desired.  
%theo 1.7  
$\frac{\partial U}{\partial t}$ and the derivative are remotely almost-periodic and ergodic for each $\nu\in(0,\infty)$.  
Thus, condition (2) is satisfied.

To prove (3). Let us choose $a=a(\nu)$, a function of $\nu$, such that $a(\nu)\rightarrow 0$ and $A(a(\nu))\xi(y,\nu)\rightarrow 0$ uniformly with respect to $y\in B_r(0)$ as $\nu\rightarrow 0$.  
Then $\nu U\rightarrow 0$ and $\nu \frac{\partial U}{\partial x}\rightarrow 0$ as $\nu \rightarrow 0$, uniformly with respect to $x\in B_r(0)$ and $t\in \mathbb{R}$, since $\nu U$ and $\nu \frac{\partial U}{\partial x}$ are bounded by $A(a(\nu))\xi(y,\nu)$.  
For every number $r<r_0$, we choose $\nu_0$ sufficiently small such that $r+a(\nu)<r_0$ for all $\nu\in(0,\nu_0)$.  
It follows from the definition of $\Delta_a(x)$ that

\begin{eqnarray*}
\int_{B_{r_{0}}} \Delta_{a(\nu)}(x-y)dy=1,\ x\in B_{r}(0),\ \nu\in(0,\nu_0).
\end{eqnarray*}
Note that
\begin{eqnarray*}
G(t,x,\nu)=\frac{\partial U}{\partial t}(t,x,\nu)	-H(t,x).
\end{eqnarray*}
Let
\begin{eqnarray*}
\overline{G}(t,x,\nu)= G(t,x,\nu)+\nu U(t,x,\nu).
\end{eqnarray*}
Since
\begin{eqnarray*}
\frac{\partial U}{\partial t}(t,x,\nu)=\int_{B_{r}(0)}\Delta_{a(\nu)}(x-y)[H(t,y)-\nu\overline{H}(t,y,\nu)]dy,
\end{eqnarray*}
we have 
\begin{eqnarray*}
\overline{G}(t,x,\nu)=\int_{B_{r}(0)}\Delta_{a(\nu)}(x-y)[H(t,y)-H(t,x)]dy.
\end{eqnarray*}
By the mean value theorem,
\begin{eqnarray*}
\Vert\overline{G}(t,x,\nu) \Vert &\leq& \sup_{0\leq\Vert x-y\Vert\leq a(\nu)} \Vert H(t,y)-H(t,x) \Vert \\
&=& \sup_{0\leq\Vert x-y\Vert\leq a(\nu)} \left\Vert \frac{\partial H}{\partial x}(t, \theta_{y} (y-x)) \right\Vert \Vert y-x\Vert,
\end{eqnarray*}
where $\theta_y \in (0,1)$. Since $\frac{\partial H}{\partial x}$ is continuous in $x$, uniformly in $t\in\mathbb{R}$, the function $\frac{\partial H}{\partial x}$ is bounded on $\mathbb{R}\times B_{r}(0)$.  
Thus, $\overline{G}(t,x,\nu)\rightarrow 0$ as $\nu\rightarrow 0$, uniformly on $\mathbb{R}\times B_{r}(0)$. Then
\begin{eqnarray*}
G(t,x,\nu)\rightarrow 0 \mbox{  when } \nu\rightarrow 0
\end{eqnarray*}
uniformly in $\bb{R}\times B_{r}(0)$. 
We have
\begin{eqnarray*}
\frac{\partial\overline{G}}{\partial x}(t,x,\nu)=\int_{B_{r_0}} \Delta_{a(\nu)}(x-y)[\frac{\partial H}{\partial x}(t,y)-\frac{\partial H}{\partial x}(t,x)]dy,
\end{eqnarray*}
using the argument employed for $\overline{G}$, it follows that
\begin{eqnarray*}
\frac{\partial\overline{G}}{\partial x}(t,x,\nu)\rightarrow 0 \mbox{  when } \nu\rightarrow 0
\end{eqnarray*}
uniformly with respecto to $(t,x)\in\bb{R}\times B_{r}(0)$. Then,
\begin{eqnarray*}
\frac{\partial G}{\partial x}(t,x,\nu)\rightarrow 0 \mbox{  when } \nu\rightarrow 0
\end{eqnarray*}
uniformly with respect to $(t,x)\in\bb{R}\times B_{r}(0)$. This proves $(3)$.

Since the four functions in $(3)$ converge to zero as $\nu\rightarrow 0$, uniformly with respect to $(t,x)\in\mathbb{R}\times B_{r}(0)$, we can define them as zero for all $(t,x)\in\mathbb{R}\times B_{r}(0)$ when $\nu=0$. Given $\nu_1>0$, let 
\begin{eqnarray*}
\Omega_1 = \{ x\ :\  x=y+\nu U(t,y,\nu),\ (t,y,\nu) \in \bb{R} \times B_{r}(0) \times [0,\nu_1] \}
\end{eqnarray*}
a compact subset of $\bb{C}^n$. Note that $\nu \frac{\partial U }{\partial y}\rightarrow 0$ as $\nu\rightarrow 0$, uniformly with respect to $(t,y)\in\mathbb{R}\times B_{r}(0)$.

Let us choose $\nu_2>0$ such that $I+\nu\frac{\partial U (t,y,\nu)}{\partial y}$ has a bounded inverse for $(t,y,\nu)\in\mathbb{R}\times B_{a(\nu_2)}(0)\times[0,\nu_2]$.  
Then, by the inverse function theorem, the change of variable \eqref{cambio_variable} has at most one solution $y\in B_{a(\nu_2)}(0)$ for each $(x,t,\nu)\in \mathbb{R}\times \Omega_1 \times [0,\nu_2]$.

For any $x_{0}\in\Omega_{1}$ there exists $\nu_{3}(x_{0})>0$ such that the change of variables \eqref{cambio_variable} has a unique solution $y=y(t,x,\nu)$, defined and continuous for $\lvert y-x_{0}\rvert\le \nu_{3}(x_{0})$, $\lvert x-x_{0}\rvert\le \nu_{3}(x_{0})$, and $0\le \nu\le \nu_{3}(x_{0})$.\\
Since $\Omega_{1}$ is compact, we can choose $\nu_{4}>0$, independent of $x_{0}$, that satisfies the same properties as $\nu_{3}(x_{0})$.\\
If $\nu_{0}=\min\{\nu_{1},\nu_{2},\nu_{4}\}$, then the change of variables defines a homeomorphism. The transformation is well defined for $(t,y,\nu)\in\bb{R}\times B_{r}(0)\times[0,\nu_{0}]$.
By \eqref{cambio_variable}, we have
\begin{eqnarray*}
\frac{dx}{dt}=\frac{dy}{dt}+\nu \frac{\partial U}{\partial y}\frac{dy}{dt}+\nu\frac{\partial U}{\partial t}.
\end{eqnarray*}

Then, by \eqref{averaging_inicial} we have
\begin{eqnarray*}
\left(I+\nu \frac{\partial U}{\partial y}\right)\frac{dy}{dt}
&=&\frac{dx}{dt}-\nu\frac{\partial U}{\partial t}\\
&=& \nu f_{\nu}(t,y+\nu U(t,y,\nu))-\nu \frac{\partial U}{\partial t}\\
&=& \nu f_{0}(y)+ \nu \left[ f_{0}(t,y)-f_0(y)- \frac{\partial U}{\partial t}\right]\\
& & +\nu \left[ f(t,y+\nu U(t,y,\nu))-f_{0}(t,y)\right]\\
&=& \nu f_0(y)+\nu \tilde{f}_{\nu}(t,y)
\end{eqnarray*}
We can see that $\tilde{f}$ has the same properties as $f$ in $\mathbb{R}\times B_{r}(0)\times[0,\nu_0]$, and moreover, $f_0(t,y)=0$ for all $(t,y)\in\mathbb{R}\times B_{r}(0)$.  
Again, by point (3), we have that $\nu \frac{\partial U}{\partial x}\rightarrow 0$ as $\nu\rightarrow 0$ uniformly on $(t,x)\in \mathbb{R}\times B_{r}(0)$.  
We can choose $\nu_0$ sufficiently small such that

\begin{eqnarray*}
\left\Vert \nu \frac{\partial U}{\partial y}\right\Vert\leq \delta <1,\ \nu\in[0,\nu_0].
\end{eqnarray*}
Then,
\begin{eqnarray*}
\left[I-\left(-\nu \frac{\partial U}{\partial y} \right) \right]^{-1}=\sum_{k=0}^{\infty}\left( -\nu \frac{\partial U}{\partial y}  \right)^k.
\end{eqnarray*}
Note that $\frac{\partial U}{\partial y}$ is remotely almost-periodic, and that $[I - (-\nu \frac{\partial U}{\partial y})]^{-1}$ is also remotely almost-periodic.  
It follows that
\begin{eqnarray*}
\frac{dy}{dt}
&=& \left[I-\left(-\nu \frac{\partial U}{\partial y} \right) \right]^{-1} (\nu f_0(y)+\nu \tilde{f}_{\nu}(t,y))\\
&=& \left(I+\sum_{k=1}^{\infty}\left( -\nu \frac{\partial U}{\partial y}  \right)^k \right)(\nu f_0(y)+\nu \tilde{f}_{\nu}(t,y))\\
&=&\nu f_0(y)+\nu \tilde{\tilde{f}}_{\nu}(t,y)
\end{eqnarray*}
We can see that $\tilde{\tilde{f}}_\nu$ has the same properties as $f$, that is, it is remotely almost-periodic and $\tilde{\tilde{f}}_\nu \rightarrow 0$.

This completes the proof.
\end{proof}

\end{document}
